\section{The exact first deviation from uniformity}
\label{sec:asymptotics}

Write $P_n(M)$ for the probability in the main theorem and
$U_{a,b}(M)=q^{-ab}$ for the law of a matrix with independent uniform entries.
The exact law gives more than convergence to $U_{a,b}$: it identifies the
first correction at every matrix, including the matrices with degenerate
combined Gram form.

Put $d=a+b$ and, for an integer $u\geq 0$, define
\[
 A_q(u)=\sum_{j=0}^{u-1}q^{2j},\qquad A_q(0)=0.
\]
The number of one-dimensional relations transverse to both coordinate
summands is
\begin{equation}
 T_1=\frac{q^r-q^{t_0}-q^{t_1}+1}{q-1}.
 \label{eq:T-one}
\end{equation}
Indeed, the radical has $(q^r-1)/(q-1)$ lines, and the lines contained in
its two disjoint coordinate intersections must be removed.

\begin{theorem}[First correction and its sharp rate]
\label{thm:first-correction}
Fix $q,a,b,H_1,H_0$ and let $n$ tend to infinity. For a proposed matrix $M$,
let $s,r,t_0,t_1$ be the invariants of its combined Gram matrix as in the
main theorem. Set
\begin{equation}
 K(M)=q^{d-1}T_1-A_q(d-s)+A_q(a-s_1)+A_q(b-s_0).
 \label{eq:K-correction}
\end{equation}
If $a+b=0$, interpret the first term as zero. Then, uniformly over all
$M\in\operatorname{Mat}_{a,b}(\mathbb F_q)$,
\begin{equation}
 q^{ab}P_n(M)=1+q^{-2n}K(M)+O(q^{-4n}).
 \label{eq:likelihood-correction}
\end{equation}
Consequently, with total variation defined as half the sum of absolute
probability differences,
\begin{equation}
 \operatorname{TV}(P_n,U_{a,b})
 =\frac{q^{-2n}}{2q^{ab}}\sum_M |K(M)|+O(q^{-4n}).
 \label{eq:TV-correction}
\end{equation}
If $a,b\geq1$, the coefficient of $q^{-2n}$ in
\eqref{eq:TV-correction} is at least $(q-1)/q$, and in particular is
strictly positive. Thus the rate $q^{-2n}$ is sharp for every fixed pair
of nonempty frames. If $a=0$ or $b=0$, the observed matrix has no entries
and its distribution is exactly uniform.
\end{theorem}

\begin{proof}
Let $z=q^{-2n}$ and define
\[
 C(u,z)=\prod_{j=0}^{u-1}(1-q^{2j}z).
\]
Factoring the powers of $q$ out of the exact law gives
\begin{equation}
 q^{ab}P_n(M)=
 \frac{\displaystyle\sum_{\ell=0}^r
   T_\ell q^{d\ell-\ell(\ell+1)/2}z^\ell
   C(d-\ell-s,z)}
 {C(a-s_1,z)C(b-s_0,z)}.
 \label{eq:rational-likelihood}
\end{equation}
For sufficiently large $n$ this follows directly from the frame-count
product. In fact it remains valid at every admissible $n$: if a quotient
is too large to embed, its product includes the zero factor with $j=n$.

The denominator has constant term one. Since
$C(u,z)=1-A_q(u)z+O(z^2)$, only $\ell=0$ and $\ell=1$ contribute to
the linear coefficient. The first contributes $-A_q(d-s)$, and the
second contributes $q^{d-1}T_1$. Inverting the denominator adds
$A_q(a-s_1)+A_q(b-s_0)$, proving \eqref{eq:likelihood-correction}.
There are finitely many matrices $M$ for the fixed parameters, so the
remainder is uniform. The inequality $\bigl||x+y|-|x|\bigr|\leq|y|$
then gives \eqref{eq:TV-correction} after summation.

For sharpness, take nonzero vectors $u\in\mathbb F_q^a$ and
$v\in\mathbb F_q^b$. The scalar $u^TMv$ is the symplectic pairing of
the nonzero vector $Yu$ with $Xv$. The latter vector is uniform on
$\mathbb F_q^{2n}\setminus\{0\}$, because the symplectic group is
transitive on nonzero vectors. Therefore
\begin{equation}
 \Pr(u^TMv=\alpha)=
 \begin{cases}
 \dfrac{q^{2n-1}-1}{q^{2n}-1},&\alpha=0,\\[6pt]
 \dfrac{q^{2n-1}}{q^{2n}-1},&\alpha\ne0.
 \end{cases}
 \label{eq:scalar-pairing-law}
\end{equation}
Under $U_{a,b}$ this scalar is uniform on $\mathbb F_q$. Its total
variation distance from uniform is exactly
\begin{equation}
 \frac{q-1}{q(q^{2n}-1)}.
 \label{eq:universal-TV-lower}
\end{equation}
Taking a function of an observation cannot increase total variation.
Thus \eqref{eq:universal-TV-lower} is a lower bound for the full matrix
distance. Multiply by $q^{2n}$ and pass to the limit in
\eqref{eq:TV-correction} to obtain the stated positive lower bound
for its leading coefficient.
\end{proof}

The constants in the $O$ terms may depend on $q,a,b,H_1,H_0$.
The theorem keeps these parameters fixed; it does not assert a uniform
approximation when the frame dimensions grow with $n$.
For another description of the obstruction to exact uniformity, any
nontrivial additive character $\psi$ satisfies
\[
 \mathbb E\,\psi(u^TMv)=-\frac{1}{q^{2n}-1},\qquad
 q^{-ab}\sum_M K(M)\psi(u^TMv)=-1.
\]
Thus every nonzero rank-one additive-character test detects the same
first-order bias, regardless of the two prescribed Gram forms.

\subsection{Two symplectic planes: a closed distribution}
\label{sec:two-planes}

Take $a=b=2$ and
\[
 H_1=H_0=J_2=\begin{pmatrix}0&1\\-1&0\end{pmatrix}.
\]
No asymptotic notation is needed in this case.

\begin{theorem}[Exact two-plane likelihood]
\label{thm:two-plane-law}
For every prime power $q$, every $n\geq1$, and every
$M\in\operatorname{Mat}_{2,2}(\mathbb F_q)$,
\begin{equation}
 q^4P_n(M)=
 \begin{cases}
 \dfrac{1+(q^3-q^2)q^{-2n}}{1-q^{-2n}},&\det M=1,\\[7pt]
 \dfrac{1-q^2q^{-2n}}{1-q^{-2n}},&\det M\ne1.
 \end{cases}
 \label{eq:two-plane-likelihood}
\end{equation}
Its total variation distance from independent uniform entries is
\begin{equation}
 \operatorname{TV}(P_n,U_{2,2})
 =\frac{(q^3-q^2+1)(q^2-1)}{q^3(q^{2n}-1)}.
 \label{eq:two-plane-TV}
\end{equation}
These formulas include characteristic two and the endpoint $n=1$.
\end{theorem}

\begin{proof}
Schur complementation and $M^TJ_2M=(\det M)J_2$ give
\[
 \operatorname{rank}\Gamma_M
 =2+\operatorname{rank}\bigl((1-\det M)J_2\bigr).
\]
Both coordinate intersections of the radical are zero. If $\det M\ne1$,
then $s=2$, $r=0$, and the numerator of the exact law is $\mathcal E_n(4,2)$.
If $\det M=1$, then $s=1$, $r=2$, and
$(T_0,T_1,T_2)=(1,q+1,1)$. The numerator is then
\[
 \mathcal E_n(4,1)+(q+1)\mathcal E_n(3,1)+\mathcal E_n(2,1).
\]
Substitute the product formula and divide by $\mathcal E_n(2,1)^2$.
Writing $z=q^{-2n}$, the critical numerator of the likelihood simplifies
to
\begin{align*}
 &(1-q^2z)(1-q^4z)
 +(q+1)q^3z(1-q^2z)+q^5z^2\\
 &\hspace{35mm}=1+(q^3-q^2)z.
\end{align*}
This proves \eqref{eq:two-plane-likelihood}.

There are $q(q^2-1)$ matrices with determinant one: choose their first
column nonzero, then choose their second column from an affine line of
size $q$. The likelihood exceeds one exactly on that set, and its
excess there is $(q^3-q^2+1)/(q^{2n}-1)$. Multiplying by its uniform
mass $(q^2-1)/q^3$ gives \eqref{eq:two-plane-TV}. At $n=1$ the second
likelihood is zero and the first gives the uniform law on
$\operatorname{SL}_2(\mathbb F_q)$, as required for two symplectic bases
of the same plane.
\end{proof}

For example, at $q=2$ the likelihood is
$(1+4\cdot2^{-2n})/(1-2^{-2n})$ on the six determinant-one matrices
and $(1-4\cdot2^{-2n})/(1-2^{-2n})$ on the other ten matrices.
The exact distance is $15/[8(2^{2n}-1)]$. For $q>2$, two invertible
matrices can fall into different likelihood classes, so ordinary matrix
rank is not a sufficient statistic.
